\chapter{Proporcionalidad, velocidad y medias}\label{ch:g8:speed}

\omterm{def:g7:prop:table}{Proporcionalidad} (\cref{ch:g7:prop}) ahora se encuentra con el mundo físico:
\omterm{def:g8:speed:def}{velocidad}, caudales, conversiones de unidades --- cantidades que son
\emph{\omterm{def:g3:division:remainder}{cocientes}} de otros dos. El capítulo se cierra con \omterm{def:g8:speed:average}{ponderado
promedios}, dónde \omterm{def:g6:propdata:def}{proporcional} El pensamiento previene un error clásico.

\section{Velocidad media}

\begin{definition}[Velocidad media]\label{def:g8:speed:def}
El \emph{velocidad promedio}\index{velocidad promedio} de un viaje es el \omterm{def:g3:division:remainder}{cociente}
\[
v = \frac{d}{t}
\qquad
\left(\text{distance divided by duration}\right),
\]
en kilómetros por hora (km/h), metros por segundo (m/s), \dots\ El
la fórmula da la vuelta: $d = v \times t $ y $ t = \frac{d}{v}$.
\end{definition}

\begin{example}\label{ex:g8:speed:basic}
Un tren cubre $270$ km en $1$ h $30$ mín. Primero convierta la hora:
$1$ h $30$ mín. $= 1.5$ h. Entonces
\[
v = \frac{270}{1.5} = 180 \text{ km/h}.
\]
Distancia en $90$ kilómetros por hora durante $2$ h $20$ mín. $= \frac73$ h:
$ d = 90 \times \frac73 = 210$ km. tiempo para $35$ km en $14$ kilómetros por hora:
$ t = \frac{35}{14} = 2.5$ h $= 2$ h $30$ mín.
\end{example}

\begin{remark}[Los minutos no son decimales.]
$1$ h $30$ min es $1.5$ h, pero $1$ h $20$ min es \emph{no} $1.2$ h:
es $1 + \frac{20}{60} = \frac43 \approx 1.33$ h. Convertir siempre
minutos a un \omterm{def:g6:fractions:def}{fracción} de una hora ($60$ mín. $= 1$ h) antes de dividir.
\end{remark}

\begin{omfigure}
\begin{tikzpicture}
\begin{axis}[omaxis,
  width=8.4cm, height=5.6cm,
  xmin=0, xmax=3.4, ymin=0, ymax=270,
  xlabel={time (h)}, ylabel={distance (km)},
  xtick={1,2,3}, ytick={80,160,240}]
  \addplot[omDef, thick, domain=0:3.2] {80*x};
  \addplot[omThm, thick, domain=0:3.2] {50*x};
  \node[omDef, font=\small, rotate=36] at (axis cs:2.35,215)
    {$80$ km/h};
  \node[omThm, font=\small, rotate=25] at (axis cs:2.6,155)
    {$50$ km/h};
\end{axis}
\end{tikzpicture}

{\small En constante \omterm{def:g8:speed:def}{velocidad}, la distancia es \omterm{def:g6:propdata:def}{proporcional} al tiempo: un
recta que pasa por el origen, cuya pendiente \emph{es} el
\omterm{def:g8:speed:def}{velocidad}.}
\end{omfigure}

\begin{method}[Velocidades de conversión]\label{met:g8:speed:convert}
Para convertir km/h a m/s: $1$ kilómetros $= 1000$ m y $1$ h $= 3600$ s, entonces
\[
v \text{ km/h} = \frac{v \times 1000}{3600} \text{ m/s}
= \frac{v}{3.6} \text{ m/s}.
\]
Para convertir m/s a km/h, multiplica por $3.6$.
\end{method}

\begin{example}\label{ex:g8:speed:convert}
$90$ kilómetros por hora $= \frac{90}{3.6} = 25$ EM. Un velocista corriendo $100$ m en
$10$ se mueve en $10$ EM $= 36$ km/h.
\end{example}

\section{Cantidades cocientes}

\begin{example}[Caudal, densidad, precio por kilo.]\label{ex:g8:speed:quotients}
\omterm{def:g8:speed:def}{Velocidad} Tiene muchos primos, todos tratados de la misma manera:
\begin{itemize}
  \item se llena un grifo $48$ L en $4$ mín: caudal
        $\frac{48}{4} = 12$ L/min, por lo que llenar un $150$ L bañera toma
        $\frac{150}{12} = 12.5$ mín;
  \item $0.6$ kg de queso cuesta $9$ euros: precio
        $\frac{9}{0.6} = 15$ euros por kilo;
  \item un auto usa $6.3$ L para $90$ km: consumo
        $\frac{6.3}{90} \times 100 = 7$ litros por $100$ km.
\end{itemize}
Identificar el \omterm{def:g3:division:remainder}{cociente}, y la fórmula de tres vías
($ q = \frac ab $, $ a = qb $, $ b = \frac aq $) hace el resto.
\end{example}

\section{Promedios ponderados}

\begin{definition}[Promedio ponderado]\label{def:g8:speed:average}
cuando los valores $ v_1, v_2, \dots $ vienen con recuentos (o pesos)
$ n_1, n_2, \dots $, su \emph{ponderado
promedio}\index{ponderado
promedio}
es
\[
\bar v = \frac{n_1 v_1 + n_2 v_2 + \dots}{n_1 + n_2 + \dots} :
\]
total de los valores, dividido por el total de los pesos.
\end{definition}

\begin{example}\label{ex:g8:speed:average}
Dos grupos realizaron una prueba: grupo 1, $12$ estudiantes, promedio
$11$; grupo 2, $18$ estudiantes, promedio $14$. promedio del conjunto
clase:
\[
\bar v = \frac{12 \times 11 + 18 \times 14}{12 + 18}
= \frac{132 + 252}{30} = \frac{384}{30} = 12.8 .
\]
\emph{No} $\frac{11 + 14}{2} = 12.5$: el grupo más grande tira del
promedio hacia el suyo propio: los pesos importan.
\end{example}

\begin{example}[Velocidad media en dos piernas.]\label{ex:g8:speed:twolegs}
un ciclista pasea $30$ km en $30$ kilómetros por hora, entonces $30$ km en $15$ km/h. El
\omterm{def:g8:speed:def}{velocidad promedio} durante todo el viaje es \emph{no}
$\frac{30 + 15}{2} = 22.5$ km/h. Calcular con la definición:
\begin{enumerate}
  \item veces: $\frac{30}{30} = 1$ h, entonces $\frac{30}{15} = 2$ h;
        total $3$ h;
  \item distancia total $60$ kilómetros, entonces
        $ v = \frac{60}{3} = 20$ km/h.
\end{enumerate}
La pierna lenta dura más, por lo que pesa más: un promedio de \omterm{def:g8:speed:def}{velocidades}
debe ser ponderado por \emph{tiempo}.
\end{example}

\section{Ejercicios}

\begin{exercise}[$\star $]\label{exo:g8:speed:1}
Calcular el \omterm{def:g8:speed:def}{velocidad promedio}: $150$ km en $2$ h; $27$ km en $45$ mín;
$100$ m en $12.5$ s.
\end{exercise}

\begin{exercise}[$\star $]\label{exo:g8:speed:2}
Calcular la distancia: $2$ sombrero en $85$ kilómetros por hora; $40$ min a las $90$ km/h.
Calcular el tiempo: $315$ km en $90$ km/h (en h y min).
\end{exercise}

\begin{exercise}[$\star $]\label{exo:g8:speed:3}
Convertir: $72$ km/h en m/s; $5$ m/s en km/h; $108$ km/h en m/s.
\end{exercise}

\begin{exercise}[$\star $]\label{exo:g8:speed:4}
Un grifo entrega $15$ L/min. ¿Cuánto tiempo para llenar un $600$ ¿Tanque L? cuantos
litros en $2$ h $30$?
\end{exercise}

\begin{exercise}[$\star $]\label{exo:g8:speed:5}
Cual es la mejor compra: $1.2$ kg de manzanas para $3.30$ euros, o
$0.8$ kilos para $2.16$ euros? Compara precios por kilogramo.
\end{exercise}

\begin{exercise}[$\star $]\label{exo:g8:speed:6}
una clase de $25$ los estudiantes tienen un promedio de $12.4$ en una prueba; otro
clase de $35$ los estudiantes tienen un promedio de $10$. Calcular el promedio de
las dos clases juntas.
\end{exercise}

\begin{exercise}[$\star $]\label{exo:g8:speed:7}
El sonido viaja a aproximadamente $340$ EM. Ves un relámpago y escuchas
el trueno $6$ segundos después. ¿A qué distancia cayó el rayo?
(al más cercano $100$ metro)?
\end{exercise}

\begin{exercise}[$\star\star $]\label{exo:g8:speed:8}
un excursionista camina $2$ sombrero en $5$ km/h, reposo $30$ min, luego camina $1$ h $30$
en $4$ km/h. Calcule la distancia total y la \omterm{def:g8:speed:def}{velocidad promedio}
\emph{incluyendo el resto} (distancia total sobre tiempo total transcurrido).
\end{exercise}

\begin{exercise}[$\star\star $]\label{exo:g8:speed:9}
Calificaciones con coeficientes: un estudiante obtuvo $15$ (coeficiente $3$), $9$
(coeficiente $2$) y $12$ (coeficiente $1$). Calcular el \omterm{def:g8:speed:average}{ponderado
promedio}. ¿Qué promedio simple (no ponderado) tendría y por qué?
¿Se diferencian?
\end{exercise}

\begin{exercise}[$\star\star $]\label{exo:g8:speed:10}
un auto conduce $120$ km en $80$ km/h y luego $120$ km en $120$ km/h.
\begin{enumerate}
  \item Calcule la duración de cada tramo, luego el \omterm{def:g8:speed:def}{velocidad promedio} encima
        todo el viaje.
  \item Explica por qué la respuesta es menor que $100$ km/h, el punto medio
        de los dos \omterm{def:g8:speed:def}{velocidades}.
\end{enumerate}
\end{exercise}

\begin{exercise}[$\star\star\star $]\label{exo:g8:speed:11}
Dos pueblos son $18$ km de distancia. Anna deja el primero en $10{:}00$
caminando en $5$ km/h hacia el segundo; Boris deja el segundo en
al mismo tiempo caminando en $4$ km/h hacia la primera. ¿A qué hora
se encuentran y ¿a qué distancia del pueblo de Anna? (Juntos ellos
cerrar la brecha en $5 + 4$ km/h.)
\end{exercise}

\section{Problema: La media armónica, o por qué se arruina el viaje de regreso
el promedio}

\begin{problem}[{Problema de fin de semana --- la velocidad promedio en distancias iguales es $\frac{2 v_1 v_2}{v_1 + v_2}$, y nunca supera el punto medio de las velocidades}]
\label{pb:g8:speed:1}
\Cref{ex:g8:speed:twolegs} y \cref{exo:g8:speed:10} ambos terminaron
con la misma sorpresa: salir a la una \omterm{def:g8:speed:def}{velocidad}, regresa en
otro, y el \omterm{def:g8:speed:def}{velocidad promedio} tierras \emph{abajo} el punto medio de
los dos. Este problema encuentra la fórmula exacta detrás de la sorpresa.
--- un nuevo tipo de promedio, el \emph{armónico
significar}\index{armónico
significar} --- demuestra que la sorpresa es una
teorema, y lo lleva a su extremo lógico: una puesta al día que ningún
\omterm{def:g8:speed:def}{velocidad} en el mundo puede lograr.

\medskip
\noindent\textbf{Parte I --- The formula for a round trip.}
Un viaje cubre una distancia $ d $ en \omterm{def:g8:speed:def}{velocidad} $ v_1$, entonces lo mismo
distancia $ d $ otra vez en \omterm{def:g8:speed:def}{velocidad} $ v_2$ (todos los números positivos).
\begin{enumerate}
  \item Exprese las dos duraciones, luego la duración total, como
        \omterm{def:g6:fractions:def}{fracciones} involucrando $ d $, $ v_1$, $ v_2$.
  \item El \omterm{def:g8:speed:def}{velocidad promedio} de todo el viaje es
        $\bar v = \dfrac{2d}{\ \dfrac{d}{v_1} + \dfrac{d}{v_2}\ }$
        --- un dos pisos \omterm{def:g6:fractions:def}{fracción}
        (\cref{ex:g8:fractions:complex}). Simplifícalo y demuestra:
        \[
        \bar v = \frac{2\, v_1 v_2}{v_1 + v_2} .
        \]
  \item Verifique la fórmula con los dos cálculos citados anteriormente:
        $ v_1 = 30$, $ v_2 = 15$ kilómetros por hora
        (\cref{ex:g8:speed:twolegs}), y $ v_1 = 80$,
        $ v_2 = 120$ kilómetros por hora
        (\cref{exo:g8:speed:10}).
  \item la distancia $ d $ ha desaparecido de la fórmula. que
        ¿Qué significa eso, concretamente, para el ciclista?
  \item Ahora supongamos que el viaje pasa lo mismo
        \emph{tiempo} $ T $ en cada \omterm{def:g8:speed:def}{velocidad}. Demuestre que el \omterm{def:g8:speed:def}{velocidad promedio} es entonces $\frac{v_1 + v_2}{2}$, el punto medio simple.
        en el lenguaje de \cref{def:g8:speed:average}: por qué
        cantidad debe ser un promedio de \omterm{def:g8:speed:def}{velocidades} siempre estar ponderado,
        ¿Y qué cambia entre los dos escenarios?
\end{enumerate}

\medskip
\noindent\textbf{Parte II --- Harmonic against arithmetic.}
Para dos números positivos $ v_1$, $ v_2$, colocar
\[
H = \frac{2\, v_1 v_2}{v_1 + v_2}
\quad\text{(their \emph{armónico
significar}),}
\qquad
A = \frac{v_1 + v_2}{2}
\quad\text{(their \emph{media aritmética}).}
\]
\begin{enumerate}[resume]
  \item Calcular $ H $ y $ A $ para las parejas $(30, 15)$ y
        $(80, 120)$. ¿Cuál de los dos medios gana cada vez?
  \item Poner $ A - H $ sobre lo común \omterm{def:g4:fractions:def}{denominador} $2(v_1 + v_2)$,
        expandir el \omterm{def:g4:fractions:def}{numerador} con las identidades notables
        (\cref{pb:g8:equations:1}), y demostrar:
        \[
        A - H = \frac{(v_1 - v_2)^2}{2\,(v_1 + v_2)} .
        \]
  \item Deduce el teorema detrás de todas las sorpresas: por positivo
        \omterm{def:g8:speed:def}{velocidades}, $ H \leq A $ siempre, con igualdad exactamente cuando
        $ v_1 = v_2$. ¿Por qué un cuadrado en el \omterm{def:g4:fractions:def}{numerador} resolver
        el asunto?
  \item Demuestra la identidad elegante $ H \times A = v_1 \times v_2$:
        las dos medias se multiplican hasta el original \omterm{def:g2:mult:def}{producto}.
  \item Un coche cubre $60$ km en $40$ km/h y $60$ km en
        $60$ km/h. Predecir el \omterm{def:g8:speed:def}{velocidad promedio} con la fórmula,
        Luego confírmalo a lo largo, con tiempos y total.
        distancia.
\end{enumerate}

\medskip
\noindent\textbf{Parte III --- The impossible catch-up.}
\begin{enumerate}[resume]
  \item Un ciclista entra en $30$ carrera de km esperando promediar
        $30$ km/h. ¿Qué tiempo total puede tomarle la carrera, en
        ¿la mayoría? Ella monta la primera $15$ km en $15$ km/h: cuanto
        ¿De ese tiempo queda presupuesto?
  \item Demuestre que el promedio esperado se ha vuelto literalmente
        imposible: computarla \omterm{def:g8:speed:def}{velocidad promedio} si ella monta el
        segundo \omterm{def:g3:division:half}{medio} en $90$ km/h y explica por qué \omterm{def:g3:numbers:evenodd}{incluso} un
        arbitrariamente enorme \omterm{def:g8:speed:def}{velocidad} la deja corta de
        $30$ km/h.
  \item Ella reduce su objetivo a $20$ km/h. Qué \omterm{def:g8:speed:def}{velocidad} en el
        segundo $15$ km lo logra? Comprueba tu respuesta con el
        fórmula de media armónica.
  \item Sólo un grifo de agua fría llena una bañera $20$ mín; un grifo caliente solo
        lo llena $30$ mín. Qué \omterm{def:g6:fractions:def}{fracción} de la tina haz ambas cosas
        Los grifos juntos se llenan por minuto y ¿cuánto tardan?
        Compara tu respuesta con $ H(20, 30)$ --- ¿Qué haces?
        aviso?
  \item Un avión recorre la misma ruta de ida y vuelta: $900$ kilómetros por hora
        con el viento, $700$ km/h en contra. Calcular el
        ida y vuelta \omterm{def:g8:speed:def}{velocidad promedio}y explica en una frase por qué
        un viento, como quiera que sople, siempre \emph{alarga} una ronda
        viaje.
\end{enumerate}
\end{problem}
